\documentclass[11pt, a4paper]{article}

% ── Packages ──────────────────────────────────────────────────────────────────
\usepackage[margin=2.5cm]{geometry}
\usepackage{amsmath, amssymb, amsthm}
\usepackage{mathtools}
\usepackage{bm}
\usepackage{microtype}
\usepackage{hyperref}
\usepackage{xcolor}
\usepackage{booktabs}
\usepackage{graphicx}
\usepackage{enumitem}
\usepackage{tcolorbox}
\usepackage{listings}
\usepackage[T1]{fontenc}
\usepackage{lmodern}

% ── Colours ───────────────────────────────────────────────────────────────────
\definecolor{copper}{HTML}{8f3e22}
\definecolor{ivory}{HTML}{f7f3ea}
\definecolor{charcoal}{HTML}{17130e}

\hypersetup{
  colorlinks=true,
  linkcolor=copper,
  urlcolor=copper,
  citecolor=copper,
}

% ── Theorem environments ──────────────────────────────────────────────────────
\theoremstyle{definition}
\newtheorem{definition}{Definition}[section]
\newtheorem{proposition}{Proposition}[section]
\newtheorem{remark}{Remark}[section]

% ── Code listings ─────────────────────────────────────────────────────────────
\lstset{
  language=Python,
  basicstyle=\small\ttfamily,
  keywordstyle=\color{copper}\bfseries,
  commentstyle=\color{gray},
  frame=single,
  framesep=4pt,
  breaklines=true,
  columns=fullflexible,
  backgroundcolor=\color{ivory},
}

% ── Notation shorthands ───────────────────────────────────────────────────────
\newcommand{\bx}{\bm{x}}
\newcommand{\bepsilon}{\bm{\epsilon}}
\newcommand{\bmu}{\bm{\mu}}
\newcommand{\bI}{\bm{I}}
\newcommand{\E}{\mathbb{E}}
\newcommand{\N}{\mathcal{N}}
\newcommand{\KL}{\mathrm{KL}}
\newcommand{\abar}{\bar{\alpha}}

% ─────────────────────────────────────────────────────────────────────────────
\title{%
  \color{copper}\textbf{Diffusion Models from Scratch}\\[4pt]
  \large TLH-1 · Thursday Learning Hours · Applied Models Track
}
\author{Thursday Learning Hours}
\date{April 10, 2026}

% ─────────────────────────────────────────────────────────────────────────────
\begin{document}
\maketitle

\begin{tcolorbox}[colback=ivory, colframe=copper, title=\textbf{Session Core Thesis}]
A diffusion model is a learned reversal of a known destruction process. The forward process is
fixed physics — Gaussian noise added step by step. The reverse is learned — a neural network
that predicts ``what was here before the noise?'' Score matching gives you the training
objective. The rest is engineering.
\end{tcolorbox}

\tableofcontents
\newpage

% ═════════════════════════════════════════════════════════════════════════════
\section{Introduction}

Diffusion models have emerged as the dominant paradigm in generative modeling, powering
Stable Diffusion~\cite{rombach2022ldm}, DALL·E 3~\cite{ramesh2022dalle2},
Sora~\cite{brooks2024sora}, and molecular design tools like
RFdiffusion~\cite{watson2022rfdiffusion}. Despite their ubiquity, most practitioners use them
as black boxes.

This report provides a self-contained mathematical derivation of Denoising Diffusion
Probabilistic Models (DDPM)~\cite{ho2020ddpm}, followed by the key extensions: DDIM
deterministic sampling~\cite{song2020ddim}, classifier-free guidance~\cite{ho2021cfg},
latent diffusion~\cite{rombach2022ldm}, and consistency models~\cite{song2023consistency}.

% ═════════════════════════════════════════════════════════════════════════════
\section{The Generative Modeling Problem}

A generative model learns the data distribution $p_\mathrm{data}(\bx)$ so we can draw novel
samples from it. Three dominant approaches:

\begin{center}
\begin{tabular}{llll}
\toprule
Model & Density & Training & Quality \\
\midrule
GANs & Implicit & Adversarial (unstable) & High \\
VAEs & Explicit (ELBO) & Stable & Moderate \\
\textbf{Diffusion} & \textbf{Iterative ELBO} & \textbf{Stable} & \textbf{State-of-the-art} \\
\bottomrule
\end{tabular}
\end{center}

% ═════════════════════════════════════════════════════════════════════════════
\section{The Forward Process}

\subsection{Markovian Gaussian Chain}

The forward process $q$ adds Gaussian noise over $T$ steps:

\begin{equation}
  q(\bx_t \mid \bx_{t-1}) = \N\!\left(\bx_t;\; \sqrt{1-\beta_t}\,\bx_{t-1},\; \beta_t\bI\right)
  \label{eq:forward_step}
\end{equation}

where $0 < \beta_1 < \beta_2 < \cdots < \beta_T < 1$ is the \textbf{noise schedule}.

\subsection{Closed-Form Marginal}

\begin{proposition}
Define $\alpha_t = 1-\beta_t$ and $\abar_t = \prod_{s=1}^t \alpha_s$. Then:
\begin{equation}
  \boxed{q(\bx_t \mid \bx_0) = \N\!\left(\bx_t;\; \sqrt{\abar_t}\,\bx_0,\;
    (1-\abar_t)\bI\right)}
  \label{eq:closed_form_marginal}
\end{equation}
\end{proposition}

\begin{proof}
By induction. For $t=1$: $q(\bx_1|\bx_0) = \N(\sqrt{\alpha_1}\bx_0, \beta_1\bI) =
\N(\sqrt{\abar_1}\bx_0, (1-\abar_1)\bI)$. Assuming the result holds for $t-1$, use the
sum-of-Gaussians identity: if $\bx_{t-1} \sim \N(\sqrt{\abar_{t-1}}\bx_0, (1-\abar_{t-1})\bI)$
and $\bx_t | \bx_{t-1} \sim \N(\sqrt{\alpha_t}\bx_{t-1}, \beta_t\bI)$, then
$\bx_t | \bx_0 \sim \N(\sqrt{\alpha_t \abar_{t-1}}\bx_0, (1 - \alpha_t\abar_{t-1})\bI)
= \N(\sqrt{\abar_t}\bx_0, (1-\abar_t)\bI)$.
\end{proof}

\textbf{Key insight:} Equation~\eqref{eq:closed_form_marginal} allows us to sample $\bx_t$
from $\bx_0$ in a single step during training:
\begin{equation}
  \bx_t = \sqrt{\abar_t}\,\bx_0 + \sqrt{1-\abar_t}\,\bepsilon, \quad \bepsilon \sim \N(\bm{0}, \bI)
\end{equation}

\subsection{Noise Schedules}

\textbf{Linear schedule} (Ho et al.): $\beta_t = \beta_\mathrm{start} +
\frac{t}{T}(\beta_\mathrm{end} - \beta_\mathrm{start})$, with $\beta_\mathrm{start} = 10^{-4}$
and $\beta_\mathrm{end} = 0.02$.

\textbf{Cosine schedule} (Nichol \& Dhariwal~\cite{nichol2021improved}):
\begin{equation}
  \abar_t = \cos\!\left(\frac{t/T + s}{1+s} \cdot \frac{\pi}{2}\right)^2 \bigg/
  \cos\!\left(\frac{s}{1+s}\cdot\frac{\pi}{2}\right)^2, \quad s = 0.008
\end{equation}
The cosine schedule distributes noise more uniformly across timesteps.

% ═════════════════════════════════════════════════════════════════════════════
\section{The Reverse Process}

\subsection{Posterior}

The forward posterior $q(\bx_{t-1}|\bx_t, \bx_0)$ is tractable:

\begin{equation}
  q(\bx_{t-1}\mid\bx_t, \bx_0) = \N\!\left(\bx_{t-1};\; \tilde{\bmu}_t(\bx_t, \bx_0),\;
  \tilde{\beta}_t\bI\right)
\end{equation}

where:
\begin{align}
  \tilde{\bmu}_t(\bx_t, \bx_0) &= \frac{\sqrt{\abar_{t-1}}\,\beta_t}{1-\abar_t}\bx_0 +
  \frac{\sqrt{\alpha_t}(1-\abar_{t-1})}{1-\abar_t}\bx_t \\
  \tilde{\beta}_t &= \frac{1-\abar_{t-1}}{1-\abar_t}\cdot\beta_t
\end{align}

\subsection{Learned Reverse}

We parameterise the reverse process as:
\begin{equation}
  p_\theta(\bx_{t-1}\mid\bx_t) = \N(\bx_{t-1};\; \bmu_\theta(\bx_t, t),\; \sigma_t^2\bI)
\end{equation}

% ═════════════════════════════════════════════════════════════════════════════
\section{The Training Objective}

\subsection{ELBO}

Maximising the variational lower bound:
\begin{equation}
  \mathcal{L} = \E_q\!\left[-\log p_\theta(\bx_T) + \sum_{t=2}^T
  \KL\!\left(q(\bx_{t-1}|\bx_t,\bx_0)\,\|\,p_\theta(\bx_{t-1}|\bx_t)\right)
  - \log p_\theta(\bx_0|\bx_1)\right]
\end{equation}

\subsection{The Ho et al. Simplification}

Parameterise the mean as a noise predictor: $\bmu_\theta(\bx_t, t) = \frac{1}{\sqrt{\alpha_t}}
\!\left(\bx_t - \frac{\beta_t}{\sqrt{1-\abar_t}}\bepsilon_\theta(\bx_t, t)\right)$.

Then each KL term simplifies to a re-weighted squared error on the noise prediction, and
Ho et al.\ show the unweighted version works best in practice:

\begin{equation}
  \boxed{\mathcal{L}_\mathrm{simple} = \E_{t, \bx_0, \bepsilon}\!\left[
  \left\|\bepsilon - \bepsilon_\theta\!\left(\sqrt{\abar_t}\bx_0 + \sqrt{1-\abar_t}\bepsilon,\; t\right)\right\|^2
  \right]}
  \label{eq:simple_loss}
\end{equation}

\textbf{Algorithm (Training):}
\begin{enumerate}[noitemsep]
  \item Sample $\bx_0 \sim q_\mathrm{data}$, $t \sim \mathrm{Uniform}(1,\ldots,T)$,
        $\bepsilon \sim \N(\bm{0},\bI)$
  \item Compute $\bx_t = \sqrt{\abar_t}\bx_0 + \sqrt{1-\abar_t}\bepsilon$
  \item Take gradient step on $\|\bepsilon - \bepsilon_\theta(\bx_t, t)\|^2$
\end{enumerate}

% ═════════════════════════════════════════════════════════════════════════════
\section{Score Matching Connection}

\begin{proposition}
The noise predictor $\bepsilon_\theta$ and the score function $\nabla_\bx \log p_t(\bx)$
are related by:
\begin{equation}
  s_\theta(\bx_t, t) = \nabla_{\bx_t} \log p_t(\bx_t) \approx
  -\frac{\bepsilon_\theta(\bx_t, t)}{\sqrt{1-\abar_t}}
\end{equation}
\end{proposition}

The denoising MSE objective~\eqref{eq:simple_loss} is equivalent to score matching
(Hyvärinen, 2005) at noise level $\sqrt{1-\abar_t}$. Training a diffusion model is therefore
training a neural network to estimate the score of the data distribution at every noise level.

% ═════════════════════════════════════════════════════════════════════════════
\section{DDPM Sampling}

\textbf{Algorithm (Sampling, 1000 steps):}
\begin{enumerate}[noitemsep]
  \item $\bx_T \sim \N(\bm{0}, \bI)$
  \item For $t = T, T-1, \ldots, 1$:
  \begin{itemize}[noitemsep]
    \item $\bmu_\theta = \frac{1}{\sqrt{\alpha_t}}\!\left(\bx_t -
          \frac{\beta_t}{\sqrt{1-\abar_t}}\bepsilon_\theta(\bx_t, t)\right)$
    \item $\bx_{t-1} = \bmu_\theta + \sqrt{\beta_t}\,\bm{z}$ where $\bm{z} \sim \N(\bm{0},\bI)$
          (or $\bm{z}=\bm{0}$ at $t=1$)
  \end{itemize}
  \item Return $\bx_0$
\end{enumerate}

% ═════════════════════════════════════════════════════════════════════════════
\section{DDIM: Deterministic Sampling}

Song et al.~\cite{song2020ddim} observe that the training objective depends only on
$q(\bx_t|\bx_0)$, not on the full Markov chain. They define a non-Markovian forward process
with the same marginals, leading to a probability-flow ODE:

\begin{equation}
  \bx_{t-1} = \sqrt{\abar_{t-1}} \cdot \underbrace{\frac{\bx_t - \sqrt{1-\abar_t}\,\bepsilon_\theta(\bx_t, t)}{\sqrt{\abar_t}}}_{\hat{\bx}_0(\bx_t)} + \sqrt{1-\abar_{t-1}} \cdot \bepsilon_\theta(\bx_t, t)
  \label{eq:ddim}
\end{equation}

This update is \textbf{deterministic} (no stochastic $\bm{z}$) and can be applied on any
subset $S \subset \{1,\ldots,T\}$ of timesteps. Using $|S|=50$ steps gives virtually
identical sample quality at $20\times$ the speed.

% ═════════════════════════════════════════════════════════════════════════════
\section{Classifier-Free Guidance}

Ho \& Salimans~\cite{ho2021cfg} jointly train conditional and unconditional models by randomly
setting the condition $y = \varnothing$ with probability $p$ during training. At sampling time:

\begin{equation}
  \tilde{\bepsilon}_\theta(\bx_t, y) = (1+w)\,\bepsilon_\theta(\bx_t, y) -
  w\,\bepsilon_\theta(\bx_t, \varnothing)
\end{equation}

$w=0$: unconditional. $w=7.5$: Stable Diffusion default — high prompt fidelity.
Geometrically, this extrapolates the score in the direction of increasing $\log p(y|\bx)$.

% ═════════════════════════════════════════════════════════════════════════════
\section{Latent Diffusion Models}

Rombach et al.~\cite{rombach2022ldm} run diffusion in the latent space of a pre-trained VAE:
\begin{equation}
  \bx \xrightarrow{E} \bm{z} = E(\bx) \quad (\text{4}{\times}\text{ smaller spatial dims})
  \xrightarrow{\mathrm{DDIM}} \tilde{\bm{z}} \xrightarrow{D} \tilde{\bx}
\end{equation}

Text conditioning is achieved via cross-attention to CLIP embeddings at every U-Net resolution
level. This enables 512$\times$512 generation on consumer GPUs by reducing the diffusion
computation from pixel space ($512^2$) to latent space ($64^2$).

% ═════════════════════════════════════════════════════════════════════════════
\section{Consistency Models}

Song et al.~\cite{song2023consistency} train a function $f_\theta(\bx_t, t)$ that directly
predicts $\bx_0$ from any point $(\bx_t, t)$ on the same ODE trajectory:

\begin{equation}
  \mathcal{L}_\mathrm{CT} = \E\!\left[d\!\left(f_\theta(\bx_{t_{n+1}}, t_{n+1}),\;
  f_{\hat\theta}(\bx_{t_n}^\Phi, t_n)\right)\right]
\end{equation}

where $\bx_{t_n}^\Phi$ is one DDIM step from $\bx_{t_{n+1}}$ and $\hat\theta$ is an EMA of
$\theta$. The result: a single forward pass generates an image — 100--1000$\times$ faster
than DDPM.

% ═════════════════════════════════════════════════════════════════════════════
\section{The U-Net Architecture}

\subsection{Time Conditioning}

Time is embedded via sinusoidal encoding (identical to the Transformer positional encoding):
\begin{equation}
  \mathrm{PE}(t, 2i) = \sin(t / 10000^{2i/d}), \quad
  \mathrm{PE}(t, 2i+1) = \cos(t / 10000^{2i/d})
\end{equation}

The embedding is projected via an MLP and added to each ResNet block's feature map via
scale+shift (affine conditioning).

\subsection{Architecture Overview}

\begin{center}
\begin{tabular}{ll}
\toprule
Component & Description \\
\midrule
Input & $\bx_t \in \mathbb{R}^{H \times W \times C}$ + timestep $t$ \\
DownBlocks & Conv + GroupNorm + SiLU + ResNet, stride 2 \\
Bottleneck & ResNet + self-attention \\
UpBlocks & Transpose Conv + skip connections + ResNet \\
Time injection & Linear $\to$ shift at each ResBlock \\
Output & $\hat\bepsilon \in \mathbb{R}^{H \times W \times C}$ \\
\bottomrule
\end{tabular}
\end{center}

\subsection{MNIST-Scale Configuration}

\begin{lstlisting}[caption={Hyperparameters for MNIST demo (runs on CPU in ~5 minutes)}]
T          = 1000    # diffusion steps
beta_start = 1e-4
beta_end   = 0.02
img_size   = 28      # MNIST
batch_size = 128
lr         = 2e-4
epochs     = 10
model_dim  = 64      # U-Net base channels
\end{lstlisting}

% ═════════════════════════════════════════════════════════════════════════════
\section{Key Takeaways}

\begin{enumerate}
  \item \textbf{The forward process is fixed — only the reverse is learned.} The forward
        destruction is analytic; only the neural network (reverse) has parameters.

  \item \textbf{The denoising loss = score matching.} These are the same objective from
        different perspectives. Predicting $\bepsilon$ is equivalent to estimating
        $\nabla_\bx \log p_t(\bx)$.

  \item \textbf{DDIM shows the reverse process is an ODE.} Stochasticity in DDPM is optional
        — the deterministic ODE view unlocks step-skipping and $20\times$ speedup.

  \item \textbf{CFG is the mechanism of control.} Extrapolating between conditional and
        unconditional scores gives prompt-following behaviour in all modern text-to-image models.

  \item \textbf{Latent diffusion = scale.} Compressing pixels into a VAE latent before
        diffusion is how generation at 512$\times$512+ resolution becomes tractable.
\end{enumerate}

% ═════════════════════════════════════════════════════════════════════════════
\section{What's Next — TLH-2: Flow Matching}

Flow matching (Lipman et al.~\cite{lipman2022flow}; Liu et al.~\cite{liu2022rectified}) asks:
what if the forward process were a \emph{straight line} rather than a Gaussian diffusion?

\begin{itemize}[noitemsep]
  \item Forward: linear interpolation $\bx_t = (1-t)\bx_0 + t\bepsilon$
  \item ODE: $\dot\bx_t = v_\theta(\bx_t, t)$ — predict the velocity field directly
  \item Loss: $\E\|\bepsilon - \bx_0 - v_\theta(\bx_t, t)\|^2$ — simpler than DDPM ELBO
  \item Straighter trajectories → fewer integration steps → faster sampling
  \item Powers Stable Diffusion 3 and FLUX
\end{itemize}

% ═════════════════════════════════════════════════════════════════════════════
\section*{References}

\begingroup
\renewcommand{\section}[2]{}
\begin{thebibliography}{99}

\bibitem{ho2020ddpm}
  Ho, J., Jain, A., \& Abbeel, P. (2020).
  Denoising Diffusion Probabilistic Models.
  \textit{NeurIPS 2020}. arXiv:2006.11239.

\bibitem{song2020ddim}
  Song, J., Meng, C., \& Ermon, S. (2020).
  Denoising Diffusion Implicit Models.
  \textit{ICLR 2021}. arXiv:2010.02502.

\bibitem{song2021sde}
  Song, Y., et al. (2021).
  Score-Based Generative Modeling through Stochastic Differential Equations.
  \textit{ICLR 2021}. arXiv:2011.13456.

\bibitem{nichol2021improved}
  Nichol, A., \& Dhariwal, P. (2021).
  Improved Denoising Diffusion Probabilistic Models.
  \textit{ICML 2021}. arXiv:2102.09672.

\bibitem{dhariwal2021diffusion}
  Dhariwal, P., \& Nichol, A. (2021).
  Diffusion Models Beat GANs on Image Synthesis.
  \textit{NeurIPS 2021}. arXiv:2105.05233.

\bibitem{ho2021cfg}
  Ho, J., \& Salimans, T. (2021).
  Classifier-Free Diffusion Guidance.
  \textit{NeurIPS Workshop 2021}. arXiv:2207.12598.

\bibitem{rombach2022ldm}
  Rombach, R., et al. (2022).
  High-Resolution Image Synthesis with Latent Diffusion Models.
  \textit{CVPR 2022}. arXiv:2112.10752.

\bibitem{ramesh2022dalle2}
  Ramesh, A., et al. (2022).
  Hierarchical Text-Conditional Image Generation with CLIP Latents (DALL·E 2).
  arXiv:2204.06125.

\bibitem{song2023consistency}
  Song, Y., et al. (2023).
  Consistency Models.
  \textit{ICML 2023}. arXiv:2303.01469.

\bibitem{peebles2023dit}
  Peebles, W., \& Xie, S. (2023).
  Scalable Diffusion Models with Transformers (DiT).
  \textit{ICCV 2023}. arXiv:2212.09748.

\bibitem{brooks2024sora}
  Brooks, T., et al. (2024).
  Video Generation Models as World Simulators (Sora).
  OpenAI Technical Report.

\bibitem{watson2022rfdiffusion}
  Watson, J. L., et al. (2022).
  De novo design of protein structure and function with RFdiffusion.
  \textit{Nature}, 620, 1089--1100.

\bibitem{lipman2022flow}
  Lipman, Y., et al. (2022).
  Flow Matching for Generative Modeling.
  \textit{ICLR 2023}. arXiv:2210.02747.

\bibitem{liu2022rectified}
  Liu, X., et al. (2022).
  Flow Straight and Fast: Learning to Generate and Transfer Data with Rectified Flow.
  \textit{ICLR 2023}. arXiv:2209.03003.

\end{thebibliography}
\endgroup

\end{document}
